\setchapterpreamble[u]{\margintoc}
\chapter{Некоторые сведения из теории графов}
\label{chpt:GraphTheoryIntro}
\tikzsetfigurename{GraphTheoryIntro_}

В данном разделе мы дадим определения базовым понятиям из теории графов, рассмотрим несколько классических задач из области анализа графов и алгоритмы их решения.
Кроме этого, поговорим о связи между линейной алгеброй и некоторыми задачами анализа графов.
Всё это понадобится нам при последующей работе.
\mytodo{Разобраться с \textbackslash BbbB, \textbackslash BbbM, \textbackslash BbbG, \textbackslash mbfscrG. Придумать и примернить общую систему стилей для имён.}

\section{Основные определения}

\begin{definition}[Помеченный ориентированный граф]
    \emph{Помеченный ориентированный граф} $\mbfscrG = \langle V, E, L \rangle$, где $V$~--- конечное множество вершин, $E$~--- конечное множество рёбер, т.ч. $E \subseteq V \times L \times V$, $L$~--- конечное множество меток на рёбрах.
    В некоторых случаях метки называют \emph{весами}%
    \sidenote{Весами метки называют, как правило, тогда, когда они берутся из какого-либо числового множества, например $\BbbR$ или $\BbbN$.}
    и тогда говорят о \emph{взвешенном} графе.
\end{definition}

\begin{definition}[Помеченный неориентированный граф]
    В случае, если для любого ребра $(u, l, v)$ в графе также содержится ребро $(v, l, u)$, говорят, что граф \emph{неориентированный}.
\end{definition}

\begin{example}\label{exmpl:garph_visualization}
    Пусть дан граф
    \begin{align*}
        \mbfscrG_1 = \langle V & =\{0, 1, 2, 3\},                    \\
        E                      & =\{(0, a, 1), (1, a, 2), (2, a, 0), \\
                               & \ \ \ (2, b, 3), (3, b, 2)\},       \\
        L                      & =\{a, b\} \rangle.
    \end{align*}
    Графическое представление графа $\mbfscrG_1$ представлено на рисунке~\ref{gr:garph_visualization}.
    \begin{marginfigure}
        \begin{center}
            \resizebox{\marginparwidth}{!}{\input{figures/edge_labelled_graph_a_3_loop_b_2_loop}}
        \end{center}
        \caption{Визуальное представление графа из примера~\ref{exmpl:garph_visualization}}
        \label{gr:garph_visualization}
    \end{marginfigure}

    Здесь $(0, a, 1)$ и $(3,b,2)$~--- это рёбра графа $\mbfscrG_1$.
    При этом $(3, b, 2)$ и $(2, b, 3)$~--- это разные рёбра.
\end{example}

В дальнейшем речь будет идти о конечных ориентированных помеченных графах.
Мы будем использовать термин \emph{граф} подразумевая именно конечный ориентированный помеченный граф, если только не оговорено противное.

Также мы будем считать, что множество вершин $V$ проиндексировано\sidenote{Более того, часто мы будем считать, что множество врешин~--- это просто отрезок $[0 \rng |V|-1]$} (см. определение~\ref{def:SetIndexing}).

\begin{definition}[Путь]
    \label{def:path}
    \emph{Путём} $\pi$ в графе $\mbfscrG$ будем называть непустую последовательность рёбер такую, что для любых двух последовательных рёбер $e_1 = (u_1, l_1, v_1)$ и $e_2 = (u_2, l_2, v_2)$ в этой последовательности, конечная вершина первого ребра является начальной вершиной второго, то есть $v_1 = u_2$.
    Будем обозначать путь из вершины $v_0$ в вершину $v_n$ как $v_0 \pi v_n$.
    Иными совами,
    \[v_0 \pi v_n = e_0,e_1, \dots, e_{n-1} = (v_0, l_0, v_1),(v_1,l_1,v_2),\dots,(v_{n-1},l_{n-1},v_n).\]
\end{definition}
Часто для представления пути мы буем использовать следующие нотации:
\begin{center}
    \input{part_01_Prep/chapter_03_GraphTheoryIntro/figures/path0.tex}
\end{center}
или
\[
    v_0 \xrightarrow[]{l_0} v_1 \xrightarrow[]{l_1} v_2 \xrightarrow[]{l_2} \ldots \xrightarrow[]{l_{n-2}} v_{n-1} \xrightarrow[]{l_{n-1}} v_n.
\]

\begin{example}
    $(0,a,1), (1,a,2) = 0 \pi_1 2$~--- путь из вершины 0 в вершину 2 в графе $\mbfscrG_1$.
    При этом $(0, a, 1), (1, a, 2), (2, b, 3), (3, b, 2) = 0 \pi_2 2$~--- это тоже путь из вершины 0 в вершину 2 в графе $\mbfscrG_1$, но он не равен $0 \pi_1 2$.
\end{example}

Кроме того, нам потребуется отношение, отражающее факт существования пути между двумя вершинами.

\begin{definition}[Достижимость в графе]
    \label{def:reach}
    Пусть дан граф $\mbfscrG = \langle V, E, L \rangle$. \emph{Отношение достижимости} в графе $\mbfscrG$~--- это отношение $P \subseteq V \times V$, такое, что $(v_i,v_j) \in P \iff \exists v_i \pi v_j$.
\end{definition}

Отметим, что в некоторых задачах удобно считать по умолчанию, что $P$ является рефлексивным (то есть $(v_i,v_i) \in P$), однако наше определение такого не допускает\sidenote{Так как путь~--- \emph{непустая} последовательность рёбер~(см. \ref{def:path}).}.
Исправить ситуацию можно явно добавив петли $(v_i,l,v_i)$ для всех вершин.

Один из способов задать граф~--- это задать его \emph{матрицу смежности}.
Неформально говоря, матрица смежности графа $\mbfscrG = \langle V, E, L \rangle$~--- это квадратная матрица $M$ размера $n \times n$, где $|V| = n$,
такая, что ячейка $M[i,j]$ хранит информацию о рёбрах из вершины $v_i$ в вершину $v_j$.
Далее мы приведём несколько классических примеров таких матриц, после чего обобщим понятие матрицы смежности удобным для дальнейшего изложения образом.

\begin{marginfigure}
    \begin{center}
        \resizebox{\marginparwidth}{!}{\input{part_01_Prep/chapter_03_GraphTheoryIntro/figures/graph1.tex}}
    \end{center}
    \caption{Неориентированный граф}
    \label{gr:graph1}
\end{marginfigure}
\begin{example}[Пример матрицы смежности неориентированного графа]
    \label{exmpl:undirectedGraphMatrix}
    Пусть дан неориентированный граф $\mbfscrG = \langle V, E\rangle$ без каких-либо меток на рёбрах (см. рисунок \ref{gr:graph1}).
    Его матрица смежности~--- булева матрица, в ячейке которой хранится $1$ только если соответствующие вершины соединены ребром:
    \[
        \begin{pmatrix}
            0 & 1 & 1 & 0 \\
            1 & 0 & 1 & 0 \\
            1 & 1 & 0 & 1 \\
            0 & 0 & 1 & 0
        \end{pmatrix}
    \]
    То есть $M[i,j] = 1 \iff (v_i,v_j) \in E$.
    Таким образом, единственная информация о ребрах~--- это факт их наличия, так как никакой дополнительной информации (меток) у рёбер нету.
    Заметим, что матрица смежности неориентированного графа всегда симметрична относительно главной диагонали.
\end{example}

\begin{marginfigure}
    \begin{center}
        \resizebox{\marginparwidth}{!}{\input{part_01_Prep/chapter_03_GraphTheoryIntro/figures/graph2.tex}}
    \end{center}
    \caption{Ориентированный граф}
    \label{gr:graph2}
\end{marginfigure}
\begin{example}[Пример матрицы смежности ориентированного графа]
    \label{example:diGraph}
    Дан ориентированный граф (см. рисунок \ref{gr:graph2}).
    Построить его булеву матрицу смежности можно применив рассуждения из предыдущего примера (\ref{exmpl:undirectedGraphMatrix}) и выглядеть она будет следующим образом:
    \[
        \begin{pmatrix}
            0 & 1 & 0 & 0 \\
            0 & 0 & 1 & 0 \\
            1 & 0 & 0 & 1 \\
            0 & 0 & 1 & 0
        \end{pmatrix}.
    \]
\end{example}

\begin{marginfigure}
    \begin{center}
        \resizebox{\marginparwidth}{!}{\input{part_01_Prep/chapter_03_GraphTheoryIntro/figures/graph3.tex}}
    \end{center}
    \caption{Помеченный граф}
    \label{gr:graph3}
\end{marginfigure}
\begin{example}[Пример матрицы смежности помеченного графа]
    \label{exmpl:matrix_for_labelled_graph}
    Пусть дан помеченный граф (см. рисунок \ref{gr:graph3}).
    В данном случае $L = \{a,b\}$.
    Так как мы не будем в данном примере запрещать параллельные рёбра, то ячейки матрицы будут содержать множества меток: метки для всех параллельных рёбер между парой вершин.
    Соответственно,  $\varnothing$ означает, что рёбер между соответствующей парой вершин нету.

    В итоге матрица смежности исходного графа выглядит следующим образом:
    \[
        \begin{pmatrix}
            \varnothing & \{a\}       & \varnothing & \varnothing \\
            \varnothing & \varnothing & \{a\}       & \varnothing \\
            \{a\}       & \varnothing & \varnothing & \{a,b\}     \\
            \varnothing & \varnothing & \{b\}       & \varnothing
        \end{pmatrix}.
    \]
\end{example}

\begin{marginfigure}
    \begin{center}
        \resizebox{\marginparwidth}{!}{\input{part_01_Prep/chapter_03_GraphTheoryIntro/figures/graph4.tex}}
    \end{center}
    \caption{Взвешенный граф}
    \label{gr:graph4}
\end{marginfigure}
\begin{example}[Пример матрицы смежности взвешенного графа]
    \label{example:apspGraph}
    Пусть дан взвешенный граф (см. рисунок \ref{gr:graph4}).
    Будем считать, что веса берутся из $\BbbR \ (L \subseteq \BbbR$).
    Нужно придумать специальное значение, которое будет храниться в ячейках, соответствующих вершинам, между которыми нету рёбер\sidenote{
        Выбрать какое-то значение из $\BbbR$ не получится, так как в общем слуае любое такое значение может быть корректной меткой существующего ребра, так как у нас метки~--- это \emph{произвольные} значения из $\BbbR$.
        В частности, не сработает <<интуитивное>> желание сказать, что в ячейках для вершин, между которыми нет рёбер, хранится $0$.
    }.
    Часто матрица смежности взвешенного графа возникает в контексте задачи поиска кратчайших расстояний между вершинами.
    В таком случае естественно предположить, что для любой вершины $v_i$ существует петля $(v_i, 0, v_i)$, хоть она явно и не изображена.
    Далее, исходя из свойств решаемой задачи\sidenote{
        Здесь нужно обратить внимание на то, что мы конструируем матрицу опираясь на решаемую задачу.
        Для иной задачи мы часто должны будем выбрать другое специальное значение.
    }, в качестве специального значения, показывающего, что рёбер нет, возьмём $+\infty$: расстояние между несоединёнными вершинами равно бесконечности.

    В результате мы получим следующую матрицу смежности:
    \[
        \begin{pmatrix}
            0       & -1.4    & +\infty & +\infty \\
            +\infty & 0       & 2.2     & +\infty \\
            0.5     & +\infty & 0       & 1.85    \\
            +\infty & +\infty & -0.76   & 0
        \end{pmatrix}.
    \]
\end{example}

В нашем последнем примере (\ref{example:apspGraph}) мы обратили внимание на то, что детали построения матрицы смежности могут зависеть от того, для решения какой задачи её планируется использовать.
В действительности, именно так и происходит: матрица смежности~--- удобный инструмент для описания некоторых свойств графа, важных при решении той или иной задачи.
Исходя из такого взгляда мы и попробуем построить обобщённое определение матрицы смежности.

Во-первых\sidenote{
    Ещё раз обратим внимание, что кроме ячеек, хранящих значения из множества меток, есть ячейки, хранящие некое <<ничего>>, когда ребро между соответствующими вершинами отсутствует.
}, нам нужен унифицированный способ из множества меток рёбер строить множество, из элементов которого будут выбираться значения ячеек матрицы смежности.
Для этого введём множество $\Opt{L}$: для графа $\mbfscrG = \langle V, E, L \rangle$, $\Opt{L}=\{\Bbbzero\} \cup L$, где $\Bbbzero \notin L$~--- дополнительное значение, соответствующее отсутствию ребра\sidenote{
    Естественным способом реализовать такую конструкцию в языках с алгебраическими типами данных является использование типа Option (OCaml, F\#), или Maybe (Haskell), или аналогов.
    Данные типы как раз предназначены для выражения того, что что-то может присутствовать (есть ребро с меткой) либо отсутствовать (нет ребра).
}: $M[i,j] = \Bbbzero \iff (v_i, _ ,v_j) \notin E$.

Во-вторых, матрица смежности редко нужна <<сама по себе>>.
Как правило, над ней необходимо выполнять какие-либо операции в рамках решения какой-то задачи.
Более того, часто операции над элементами матрицы нужны просто для того, чтобы аккуратно определить, как именно сохраняется информация о метках параллельных рёбер\sidenote{В примере~\ref{exmpl:matrix_for_labelled_graph} мы решили эту проблему используя множества, но это не всегда уместно.}.
Дабы обеспечить эту возможность, мы будем либо сразу говорить, что матрица определена над какой-либо алгебраической структурой\sidenote{
    Конкретный тип структуры будет зависеть от решаемой задачи. Где-то хватит моноида, где-то потребуется полукольцо, а где-то ещё более сложная структура.
}, использующей в качестве носителя $\Opt{L}$, либо определять соответствующую структуру отдельно\sidenote{
    Такой подход может быть удобен с точки рения разделения данных и операций над ними.
    В частности, мы можем опеределить несколько структур с одинаковым носителем и разными операциями.
}.

Для определения обобщённой матрицы смежности нам будет удобно пользоваться первым вариантом, так как он позволяет естественным образом получить представление для параллельных рёбер\sidenote{Хотя при решении прикладных задач чаще используется второй, так как во многих языках программирования удобнее разделить стуктуру данных и функции работы с ней.}.

\begin{definition}[Матрица смежности]
    \emph{Матрица смежности} графа $\mbfscrG = \langle V, E, L \rangle$~--- это квадратная матрица $M$ размера $n \times n$, где $|V| = n$, построенная над коммутативным\sidenote{Коммутативность нам нужна для того, чтобы при построении матрицы не думать, в каком порядке добавлять параллельные рёбра.} моноидом $\BbbG = (\Opt{L}, \circ)$.
    При этом \[M[i,j] = \underset{{(i, l, j) \in E}}{\bigcirc} l\], где $\underset{\varnothing}{\bigcirc} = \Bbbzero$ и $\Bbbzero$~--- нейтральный элемент относительно $\circ$.
\end{definition}

Заметим, что наше определение матрицы смежности требует некоторой аккуратности при соотнесении с ним данных выше примеров.
Во-первых, в нашем определении присутствует моноид, который в примерах не фигурировал вовсе.
Во-вторых, выбор специального значения и способ представления информации о параллельных рёбрах в примерах выбирался исходя из некоторой интуитивной <<естественности>>, но в реальности нам необходимо аккуратно согласовывать данный выбор с решаемой задачей и соответствующими алгебраическими структурами.
Таким образом, свойства структуры $\BbbG$, а значит и детали её построения, зависят от задачи, в рамках которой рассматривается граф.
А значит, и матрица смежности конструируется исходя из задачи.

В ряде случаев граф удобно представлять не одной матрицей смежности, а набором матриц.
Рассмотрим частный случай, в котором матрица смежности строится над моноидом $\BbbG = (\mathcal{P}(L), \cup)$, где $\mathcal{P}(L)$~--- множество всех подмножеств $L$, операция $\cup$~--- объединение множеств, а $\Bbbzero = \varnothing$.
Иными словами, ячейка $M[i,j]$ содержит множество всех меток рёбер, соединяющих вершину $i$ с вершиной $j$.
Именно такую матрицу мы построили в примере~\ref{exmpl:matrix_for_labelled_graph}.

Далее, для каждого элемента $l \in L$ можно построить отдельную булеву матрицу, несущую информацию о том, существует ли в графе ребро с меткой $l$.

\begin{definition}[Булева декомпозиция матрицы смежности]
    \label{def:BoolDecomposition}
    \emph{Булевой декомпозицией} матрицы смежности $M$ графа $\mbfscrG = \langle V, E, L \rangle$ (построенной над $\mathcal{P}(L)$) называется семейство булевых матриц $\mathcal{M} = \{M_l \mid l \in L \}$, где каждая $M_l$ размера $|V| \times |V|$ определена как
    \[ M_l[i,j] = \begin{cases}
        1, & \text{если } l \in M[i,j] \\
        0, & \text{иначе}.
    \end{cases} \]
\end{definition}

Иными словами, $M_l[i,j] = 1$ тогда и только тогда, когда в исходном графе существует ребро $(i,l,j)$.

\begin{example}[Булева декомпозиция для матрицы из примера~\ref{exmpl:matrix_for_labelled_graph}]
    \label{exmpl:bool_decomposition}
    Напомним матрицу смежности $M$ из примера~\ref{exmpl:matrix_for_labelled_graph}:
    \[
        \begin{pmatrix}
            \varnothing & \{a\}       & \varnothing & \varnothing \\
            \varnothing & \varnothing & \{a\}       & \varnothing \\
            \{a\}       & \varnothing & \varnothing & \{a,b\}     \\
            \varnothing & \varnothing & \{b\}       & \varnothing
        \end{pmatrix}.
    \]

    Так как $L = \{a, b\}$, булева декомпозиция состоит из двух матриц~--- $M_a$ и $M_b$:
    \[
        M_a = \begin{pmatrix}
            0 & 1 & 0 & 0 \\
            0 & 0 & 1 & 0 \\
            1 & 0 & 0 & 1 \\
            0 & 0 & 0 & 0
        \end{pmatrix},
        \quad
        M_b = \begin{pmatrix}
            0 & 0 & 0 & 0 \\
            0 & 0 & 0 & 0 \\
            0 & 0 & 0 & 1 \\
            0 & 0 & 1 & 0
        \end{pmatrix}.
    \]
    То есть $\mathcal{M}_G = \{M_a, M_b\}$.
\end{example}

Как итог, уже можно заметить, что введение алгебраических структур как абстракции позволяет достаточно унифицированным образом смотреть на различные графы и их матрицы смежности.
Далее мы увидим, что данный путь позволит решать унифицированным образом достаточно широкий круг задач, связанных с анализом путей в графах.
Но сперва мы рассмотрим некоторые классические задачи анализа графов, которые понадобятся нам далее, и покажем их связь с линейной алгеброй.

\section{Задачи поиска путей}

Одна из классических задач анализа графов~--- это задача поиска путей между вершинами с различными ограничениями.

При этом возможны различные постановки задачи.
С одной стороны, постановки различаются тем, что именно мы хотим получить в качестве результата.
Здесь наиболее частыми являются следующие варианты.

\begin{itemize}
    \item Наличие хотя бы одного пути, удовлетворяющего ограничениям, в графе.
          В данном случае не важно, между какими вершинами существует путь, важно лишь наличие его в графе.
    \item Наличие пути, удовлетворяющего ограничениям, между некоторыми вершинами: задача достижимости.
          При данной постановке задачи, нас интересует ответ на вопрос о достижимости вершины \circled{$v_i$} из вершины \circled{$v_j$} по пути, удовлетворяющему ограничениям.
          Такая постановка требует лишь проверить существование пути, но не его предоставления в явном виде.
    \item Поиск одного пути, удовлетворяющего ограничениям: необходимо не только установить факт наличия пути, но и  предъявить его.
          При этом часто подразумевается, что возвращается любой путь, являющийся решением, без каких-либо дополнительных ограничений.
          Хотя, например, в некоторых задачах дополнительное требование простоты или наименьшей длины выглядит достаточно естественным.
    \item Поиск всех путей: необходимо предоставить все пути, удовлетворяющие заданным ограничениям.
\end{itemize}

С другой стороны, задачи могут различаться ещё и тем, как фиксируются множества стартовых и конечных вершин.
Здесь возможны следующие варианты:
\begin{itemize}
    \item от одной вершины до всех,
    \item между всеми парами вершин,
    \item между фиксированной парой вершин,
    \item между двумя множествами вершин $V_1$ и $V_2$, что подразумевает решение задачи для всех $(v_i,v_j) \in V_1 \times V_2$.
\end{itemize}

Стоит отметить, что последний вариант является самым общим и остальные~--- лишь его частные случаи.
Однако этот вариант часто выделяют отдельно, подразумевая, что остальные варианты в него не включаются\sidenote{
    Часто это связано с тем, что для некоторых частных случаев можно построить специализированный алгоритм, который будет по каким-то показателям лучше, чем алгоритм для задачи общего вида.
}.

В итоге, перебирая возможные варианты желаемого результата и способы фиксации стартовых и финальных вершин, мы можем сформулировать достаточно большое количество задач.
Например, задачу поиска всех путей между двумя заданными вершинами, задачу поиска одного пути от фиксированной стартовой вершины до каждой вершины в графе, или задачу достижимости между всеми парами вершин.

Часто поиск путей сопровождается изучением их свойств, что далее приводит к формулированию дополнительных ограничений на пути в терминах этих свойств.
Например, можно потребовать, чтобы пути были простыми или не проходили через определённые вершины.
Один из естественных способов описывать свойства и, как следствие, задавать ограничения~--- это использовать ту алгебраическую структуру, из которой берутся веса рёбер графа%
\sidenote{На самом деле здесь наблюдается некоторая двойственность.
    С одной стороны, действительно, удобно считать, что свойства описываются в терминах некоторой заданной алгебраической структуры.
    Но, вместе с этим, структура подбирается исходя из решаемой задачи.}.

Предположим, что дан граф $\mbfscrG = \langle V, E, L\rangle $, где $L = (S, \oplus, \otimes)$~--- это полукольцо.
Тогда изучение свойств путей можно описать следующим образом:
\begin{equation}
    \label{eq:algPathProblem}
    \left\{\ (v_i, v_j, c) \mid c = \bigoplus_{v_i \pi_k v_j} \bigotimes_{(u, l, v) \in \pi } l \ \right\}.
\end{equation}

Иными словами, для каждой пары вершин, для которой существует хотя бы один путь, их соединяющий, мы агрегируем (с помощью операции $\oplus$ из полукольца) информацию обо всех путях между этими вершинами.
При этом информация о пути получается как свёртка меток рёбер пути с использованием операции $\otimes$\sidenote{Заметим, что детали свёртки вдоль пути зависят от свойств полукольца (и от решаемой задачи).
    Так, если полукольцо коммутативно, то нам не обязательно соблюдать порядок рёбер.
    В дальнейшем мы увидим, что данные особенности полукольца существенно влияют на особенности алгоритмов решения соответствующих задач.}.

Естественным требованием (хотя бы для прикладных задач, решаемых таким способом) является существование и конечность указанной суммы.
На данном этапе мы не будем касаться того, какие именно свойства полукольца могут нам обеспечить данное свойство, однако в дальнейшем будем считать, что оно выполняется.
Более того, будем стараться приводить частные для конкретной задачи рассуждения, показывающие, почему это свойство выполняется в рассматриваемых в задаче ограничениях.

Описанная выше задача общего вида называется анализом свойств путей алгебраическими методами (Algebraic Path Problem)~\sidecite{Baras2010PathPI} и предоставляет общий способ для решения широкого класса прикладных задач%
\sidenote{В работе \enquote{Path Problems in Networks}~\cite{Baras2010PathPI} собран действительно большой список прикладных задач с описанием соответствующих полуколец.
    Сводная таблица на страницах 58--59 содержит 29 различных прикладных задач и соответствующих полуколец.}.
Наиболее известными являются такие задачи, как построение транзитивного замыкания графа и поиск кратчайших путей (All Pairs Shortest Path или APSP).
Далее мы подробнее обсудим эти две задачи и предложим алгоритмы их решения.


\section{Анализ путей в графе и линейная алгебра}
\label{sec:PathAlgebra}
\tikzsetfigurename{GraphTheory_PathAlgebra_}

Часто поиск путей сопровождается изучением их свойств, что далее приводит к формулированию дополнительных ограничений на пути в терминах этих свойств.
Например, можно потребовать, чтобы пути были простыми или не проходили через определённые вершины.
Один из естественных способов описывать свойства и, как следствие, задавать ограничения~--- это использовать ту алгебраическую структуру, из которой берутся веса рёбер графа%
\sidenote{На самом деле здесь наблюдается некоторая двойственность.
    С одной стороны, действительно, удобно считать, что свойства описываются в терминах некоторой заданной алгебраической структуры.
    Но, вместе с этим, структура подбирается исходя из решаемой задачи.}.

Предположим, что дан граф $\mbfscrG = \langle V, E, L\rangle $, где над $L$ построено полукольцо $\algebraic{L} = (S, \opAdd, \opMult)$\label{itm:otimesIntro}.
Тогда изучение свойств путей можно описать следующим образом:
\begin{equation}
    \label{eq:algPathProblem}
    \left\{\ (v_i, v_j, c) \mid c = \bigoplus_{v_i \pi_k v_j} \bigotimes_{(u, l, v) \in \pi } l \ \right\}.
\end{equation}

Иными словами, для каждой пары вершин, для которой существует хотя бы один путь, их соединяющий, мы агрегируем (с помощью операции $\opAdd$ из полукольца) информацию обо всех путях между этими вершинами.
При этом информация о пути получается как свёртка меток рёбер пути с использованием операции%
\sidenote{Заметим, что детали свёртки вдоль пути зависят от свойств полукольца (и от решаемой задачи).
    Так, если полукольцо коммутативно, то нам не обязательно соблюдать порядок рёбер.
    В дальнейшем мы увидим, что данные особенности полукольца существенно влияют на особенности алгоритмов решения соответствующих задач.}
$\opMult$.

Естественным требованием (хотя бы для прикладных задач, решаемых таким способом) является существование и конечность указанной суммы.
На данном этапе мы не будем касаться того, какие именно свойства полукольца могут нам обеспечить данное свойство, однако в дальнейшем будем считать, что оно выполняется.
Более того, будем стараться приводить частные для конкретной задачи рассуждения, показывающие, почему это свойство выполняется в рассматриваемых в задаче ограничениях.

Описанная выше задача общего вида называется анализом свойств путей алгебраическими методами (Algebraic Path Problem)~\sidecite{Baras2010PathPI} и предоставляет общий способ для решения широкого класса прикладных задач%
\sidenote{В работе \enquote{Path Problems in Networks}~\cite{Baras2010PathPI} собран действительно большой список прикладных задач с описанием соответствующих полуколец.
    Сводная таблица на страницах 58--59 содержит 29 различных прикладных задач и соответствующих полуколец.}.
Наиболее известными являются такие задачи, как построение транзитивного замыкания графа и поиск кратчайших путей (All Pairs Shortest Path или APSP).
Далее мы подробнее обсудим эти две задачи и предложим алгоритмы их решения.


 В данной главе мы рассмотрим некоторые связи%
 \sidenote{Связь между графами и линейной алгеброй~--- обширная область, в которой можно даже выделить отдельные направления, такие как спектральная теория графов.
     С точки зрения практики данная связь также подмечена давно и более полно с ней можно ознакомиться, например, в работах~\cite{doi:10.1137/1.9780898719918, Davis2018Algorithm9S}.
     Кроме этого, много полезной информации можно найти на сайте \url{https://graphblas.github.io/GraphBLAS-Pointers/}.}
 между графами и операциями над ними и матрицами и операциями над матрицами.

 Как мы видели, достаточно естественное представление графа~--- это его матрица смежности.
 Далее можно заметить некоторое сходство между определением матричного умножения~\ref{def:MxM} и мыслями, которыми мы руководствовались, вводя операцию $\opMult$~\ref{itm:otimesIntro}.
 Действительно, пусть есть матрица $M$ размера $n \times n$ над $\algebraic{G} = (S, \opAdd, \opMult)$~--- матрица смежности графа.
 Умножим её саму на себя: вычислим $M'= M \mmultFrom{G} M$.
 Тогда, по~\ref{def:MxM}, $M'[i, j] = \bigoplus_{l \in [0 \rng n-1]} M[i, l] \opMult M[l, j]$.
 Размер $M'$ также $n \times n$.
 То есть $M'$ задаёт такой граф, что в нём будет путь со свойством, являющимся агрегацией свойств всех путей, составленных из двух подпутей в исходном графе.
 А именно, если в исходном графе есть путь из $i$ в $l$ с некоторым свойством (его значение хранится в $M[i, l]$), и был путь из $l$ в $j$ (его значение хранится в $M[l,j]$), то в новом графе свойство $M[i, l] \opMult M[l, j]$ аддитивно (используя $\opAdd$) учтётся в свойстве пути из $i$ в $j$.

 \begin{example}
    \label{exmpl:matrixMultAndPaths}
     \NiceMatrixOptions{code-for-first-row = \color{red},
                     code-for-first-col = \color{blue},
                     code-for-last-row = \color{red},
                     code-for-last-col = \color{blue}
                     }

    Проиллюстрируем связь матричного умножения и вычисления свойств путей.
    Пусть матрица $A$ содержит информацию о рёбрах из стартовых вершин в промежуточные, а $B$~--- из промежуточных в конечные.
    Тогда $C = A \cdot B$ агрегирует информацию о путях длины~$2$, составленных из двух подпутей.
     Схематично, вычисление элемента $c_{ij}$ показано ниже ($c_{ij} = \bigoplus_{k} a_{ik} \opMult b_{kj}$):
\[
\begin{pNiceArray}{cccc}[first-col,last-col,nullify-dots]
  &  & \Cdots &  &  & \\
i & a_{i0} & \Cdots & & a_{i(k-1)} & i \\
  &  & \Cdots &  &  & \\
  &  & \Cdots &  &  &
\end{pNiceArray}
\times
\begin{pNiceArray}{cccc}[first-row,last-row=5,nullify-dots]
 & & j &                             \\
   &              & b_{0j} &         \\
  \Vdots & \Vdots & \Vdots & \Vdots  \\
   &              &        &         \\
   &              & b_{(k-1)j} &     \\
 & & j &
\end{pNiceArray}
=
\begin{pNiceArray}{cccc}[first-col,last-col,first-row,last-row=5,nullify-dots]
   &         &        & j       &        &   \\
   &         &         & \Vdots  &        &   \\
 i & \Cdots  & \Cdots  & c_{ij} & \Cdots & i \\
   &          &        & \Vdots &        &  \\
   &          &        & \Vdots &        &  \\
    &         & & j       &        &
\end{pNiceArray}
\]

    Соответствующий фрагмент графа для двух промежуточных вершин $k_0$ и $k_1$ изображён на рисунке~\ref{fig:algebraic_path_problem_path_example}.

 \begin{marginfigure}
    \begin{tikzpicture}[shorten >=1pt,auto]
    \node[state,fill=blue!20] (q_0)                      {$i$};
    \node[state] (q_1) [right = of q_0 ]    {$k_0$};

    \node[state] (q_2)  [below = of q_1 ]    {$k_1$};
    \node[state,fill=red!20] (q_3)  [right = of q_1]    {$j$};

    \path[->]
    (q_0) edge  node {$a_0$} (q_1)
    (q_0) edge[right]  node {$a_1$} (q_2)
    (q_1) edge  node {$b_0$} (q_3)
    (q_2) edge[right]  node {$b_1$} (q_3);


 \path[->,red]
   (q_0) edge[bend left=40]  node {$ a_0 \otimes b_0 \oplus a_1 \otimes b_1 $} (q_3)
   ;

\end{tikzpicture}

\caption{Вычисление свойств пути при умножении матриц смежности}
\label{fig:algebraic_path_problem_path_example}
\end{marginfigure}
 \end{example}

 Таким образом, произведение матриц смежности соответствует обработке информации о путях интересующим нас образом.
 Это наблюдение позволяет предложить решение задач анализа свойств путей, основанное на операциях над матрицами.
 Проиллюстрируем этот подход на двух классических задачах: транзитивном замыкании и поиске кратчайших путей.

 Рассмотрим сперва более простую задачу~--- транзитивное замыкание графа, которая также является частным случаем~\eqref{eq:algPathProblem}.
 Здесь метки рёбер не несут дополнительной информации: нас интересует лишь факт достижимости.
 Соответствующее полукольцо~--- булево: $\algebraic{B} = \langle \{0,1\}, \lor, \land \rangle$, где $\Bbbzero = 0$ (пути нет), $\Bbbone = 1$ (путь существует).
 Матрица смежности $M$ над $\algebraic{B}$ определяет отношение смежности вершин: $M[i,j] = 1$ тогда и только тогда, когда существует ребро из $i$ в $j$.

 Матричное умножение над $\algebraic{B}$ принимает вид
 \[
     (M \cdot N)[i,j] = \bigvee_{k \in [0 \rng n-1]} M[i,k] \land N[k,j].
 \]
 Иными словами, $(M \cdot N)[i,j] = 1$ в точности тогда, когда существует вершина~$k$, для которой $M[i,k] = N[k,j] = 1$.

 Следовательно, $M^d$ содержит информацию о путях длины ровно~$d$: $M^d[i,j] = 1$ тогда и только тогда, когда существует путь из $i$ в $j$, состоящий ровно из $d$ рёбер.


 \begin{example}
     \label{exmpl:transitiveClosure}
     Рассмотрим граф из примера~\ref{exmpl:diGraph}.

     \begin{marginfigure}
         \resizebox{\marginparwidth}{!}{\begin{tikzpicture}[]
  % Apex angle 60 degrees
  \node[isosceles triangle,
      isosceles triangle apex angle=60,
      draw=none,fill=none,
      minimum size=2cm] (T60) at (3,0){};

  \node[state] (q_0) at (T60.right corner) {$0$};
  \node[state] (q_1) at (T60.left corner)  {$1$};
  \node[state] (q_2) at (T60.apex)         {$2$};
  \node[state] (q_3) [right=1.5cm of q_2]  {$3$};
  \path[->]
    (q_0) edge (q_1)
    (q_1) edge (q_2)
    (q_2) edge (q_0)
    (q_2) edge [bend left] (q_3)
    (q_3) edge [bend left] (q_2);
  \path[->, red]
    (q_0) edge [bend left] (q_2)
    (q_1) edge [bend right] (q_0)
    (q_1) edge [bend left] (q_3)
    (q_2) edge [bend right] (q_1)
    (q_2) edge [loop above] (q_2)
    (q_3) edge [bend left] (q_0)
    (q_3) edge [loop right] (q_3);
\end{tikzpicture}
}
         \caption{Исходные рёбра графа (чёрным) и пути длины~$2$ (красным)}
         \label{fig:transitiveClosure_paths_length2}
     \end{marginfigure}

     Его булева матрица смежности~--- $M$, а квадрат этой матрицы~--- $M^2$:
     \[
         M = \begin{pmatrix}
             0 & 1 & 0 & 0 \\
             0 & 0 & 1 & 0 \\
             1 & 0 & 0 & 1 \\
             0 & 0 & 1 & 0
         \end{pmatrix},\qquad
         M^2 = \begin{pmatrix}
             0 & 0 & 1 & 0 \\
             1 & 0 & 0 & 1 \\
             0 & 1 & 1 & 0 \\
             1 & 0 & 0 & 1
         \end{pmatrix}.
     \]
     Результат этого шага представлен на рисунке~\ref{fig:transitiveClosure_paths_length2}, где красным цветом показаны рёбра, соответствующие путям длины~$2$ в исходном графе.
     Видим, что, например, $M^2[0,2] = 1$, так как существует путь $0 \to 1 \to 2$ длины~$2$.
     При этом $M^2[0,0] = 0$, поскольку путь длины ровно~$2$ из $0$ в $0$ отсутствует, а рефлексивности нет (в исходной матрице на диагонали стоят нули).
 \end{example}

 В задаче транзитивного замыкания нас интересуют пути любой длины, а не только фиксированной.
 Для этого необходимо замкнуть отношение смежности относительно транзитивности.
 Добавим петли во все вершины с весом $\Bbbone$: положим $M' = I \lor M$, где $I$~--- единичная матрица ($I[i,i] = \Bbbone$, $I[i,j] = \Bbbzero$ при $i \neq j$).
 Тогда $(M')^d$ содержит информацию о путях длины \emph{не более}~$d$~\sidenote{Работает линейная алгебра: $(M+I)^2 = M^2 + IM + MI + I^2$. То есть появляются пути длины два и остаются длины один. Далее можно продолжать в том же духе.}.
 Поскольку в ориентированном графе ни один простой путь не содержит более $n-1$ ребра, транзитивное замыкание соответствует $(M')^{n-1}$, причём суммировать промежуточные степени не нужно~--- они уже содержатся в старшей.

 Таким образом, задача транзитивного замыкания свелась к возведению булевой матрицы в степень%
 \sidenote{Рекомендуем построить транзитивное замыкание для примера~\ref{exmpl:transitiveClosure}.}
  $n-1$, что является частным случаем общего подхода~\eqref{eq:algPathProblem}.

 Заметим, что описанная схема~--- возведение матрицы смежности с петлями нейтрального веса в степень $n-1$ над подходящим полукольцом~--- является общей для широкого класса задач анализа свойств путей.
 Рефлексивность гарантирует, что $M^k$ содержит информацию о путях длины не более $k$, а ограниченность простых путей $n-1$ рёбрами~--- что степень $n-1$ достаточна.

 Применим эту схему к задаче APSP.
 Здесь используется тропическое полукольцо $\algebraic{T} = \langle \BbbR \cup \{+\infty\}, \min, + \rangle$, а рефлексивность обеспечивается нулями на главной диагонали матрицы смежности $M$.
 Тогда $M_k = M^k$~--- матрица кратчайших путей длины не более $k$, и, как и в случае транзитивного замыкания, достаточно вычислить $M_{n-1}$.
 Здесь мы также пользуемся тем, что в ориентированном графе без отрицательных циклов кратчайший путь не может содержать более $n$ рёбер%
 \sidenote{Раз отрицательных циклов нет, то проходить через одну вершину, коих $n$, больше одного раза бессмысленно.}.

 Таким образом, решение APSP сведено к произведению матриц над $\algebraic{T}$, однако вычислительная сложность наивного подхода слишком большая: $O(n \cdot \SMM{\algebraic{T}}{n})$.

 Благодаря рефлексивности%
 \sidenote{Как мы видели на примере транзитивного замыкания, матрица $M_k$ содержит информацию о путях длины не более $k$, а не ровно $k$.}%
 , для вычисления $M_{n-1}$ можно применить \emph{быстрое возведение в степень}.
 Будем последовательно возводить матрицу в квадрат, получая матрицы $M_2$, $M_4$, $M_8$ и так далее.

 В итоге получим следующую последовательность вычислений:
 \begin{gather*}
     M_1 = M \\
     M_2 = M_1 \cdot M_1 = M^2 \\
     M_4 = M_2 \cdot M_2 = M^4 \\
     M_8 = M_4 \cdot M_4 = M^8 \\
     \vdots \\
     M_{2^k} = M_{2^{k-1}} \cdot M_{2^{k-1}} = M^{2^k}
 \end{gather*}

 Здесь мы предполагаем, что $n-1 = 2^k$ для некоторого $k$.
 На практике такое верно далеко не всегда.
 Если это условие не выполняется, то необходимо взять ближайшую сверху степень двойки%
 \sidenote{
     Кажется, что это приведёт к избыточным вычислениям.
     Попробуйте оценить, на сколько много лишних вычислений будет сделано в худшем случае.}%
 , вычислить $M_{2^{\lceil \log(n-1) \rceil}}$, после чего результат будет содержать и искомую матрицу $M_{n-1}$.

 Теперь вместо $n$ итераций нам нужно $\log{n}$, а итоговая сложность решения~--- $O(\log{n} \cdot \SMM{\algebraic{T}}{n})$%
 \sidenote{
     Заметим, что это оценка для худшего случая.
     На практике при использовании данного подхода можно прекращать вычисления как только при двух последовательных шагах получились одинаковые матрицы ($M_i = M_{i-1}$).
     Это приводит нас к понятию \emph{неподвижной точки}, обсуждение которого лежит за рамками повествования.}%
 .
 Данный алгоритм%
 \sidenote{Пример решения APSP с помощью быстрого возведения в степень: \url{http://users.cecs.anu.edu.au/~Alistair.Rendell/Teaching/apac_comp3600/module4/all_pairs_shortest_paths.xhtml}.}
 называется \emph{быстрым возведением в степень} (repeated squaring).

 Итак, быстрое возведение в степень сводит APSP к $O(\log n)$ умножениям матриц над тропическим полукольцом.
 Однако, как обсуждалось в разделе~\ref{sec:TheoreticalComplexity}, быстрые алгоритмы матричного умножения (алгоритм Штрассена и его развитие, дающие оценку $O(n^\omega)$ с $\omega < 2.373$) требуют возможности вычитания, то есть работают над кольцами, но не над полукольцами.
 Для умножения матриц над полукольцами лучшие известные алгоритмы являются \emph{слегка субкубическими} (mildly subcubic): наилучшая оценка составляет $\hat{O}(n^3 / \log^4 n)$~\sidecite{10.1007/978-3-662-47672-7_89}.
 Таким образом, и APSP через этот подход получает лишь слегка субкубическую сложность $O(n^3 \log n / \log^4 n)$.

 Вопрос о существовании \emph{истинно} субкубического алгоритма для общего случая APSP остаётся открытым~\sidecite{Chan2010} и является одной из центральных проблем мелкозернистой сложности.
 Вместе с тем активно развиваются подходы, дающие истинно субкубические алгоритмы для структурированных частных случаев Min-Plus произведения и соответствующих подклассов APSP: Вильямс и Сюй (SODA'20)~\sidecite{3381089.3381091} предложили алгоритм для матриц, блоки которых близки к матрицам константного ранга; Сюй (2021)~\sidecite{xu_2021}~--- для матриц с константным аддитивным приближённым рангом и монотонных матриц; Чи, Дуан и Се (SODA'22)~\sidecite{doi:10.1137/1.9781611977073.60}~--- для ограниченно-разностных матриц с оценкой $O(n^{2.779})$.

 Заметим, что быстрое возведение в степень, использованное нами для транзитивного замыкания и APSP, возможно только при ассоциативности произведения матриц\sidenote{Ассоциативность гарантируется свойствами полукольца.}.
 Для более бедных алгебраических структур, не обеспечивающих ассоциативность, данный приём неприменим.

\mycomment{Резкий переход к разделу про GraphBLAS и программные абстракции. Нужен связующий абзац-мостик.}

\section{Прикладные особенности}
\label{sec:AppliedAspects}
\tikzsetfigurename{GraphTheory_Applied_}

В предыдущих разделах мы определили полугруппы, моноиды, группы, полукольца и кольца как <<честные>> алгебраические структуры: сформулировали аксиомы, которым они должны удовлетворять, и привели примеры.
На практике же при реализации алгоритмов анализа графов эти аксиомы часто соблюдаются лишь частично, а на первый план выходят вопросы представления данных, эффективности операций и удобства программного интерфейса.
Данный раздел посвящён именно таким прикладным аспектам~\sidecite{7761646}.

Начнём с простого наблюдения.
С точки зрения программиста естественно уметь выполнять матричное умножение в ситуации, когда аргументы и результат имеют разный тип.
Пусть даны матрицы
\[
    A: \texttt{Matrix}\langle T_1 \rangle, \qquad
    B: \texttt{Matrix}\langle T_2 \rangle,
\]
а результирующая матрица должна иметь тип
\[
    C: \texttt{Matrix}\langle T_3 \rangle.
\]
Для реализации такого обобщённого умножения потребуется определить всего три сущности:
\[
    \otimes: T_1 \times T_2 \to T_3, \qquad
    \oplus: T_3 \times T_3 \to T_3, \qquad
    \Bbbzero \in T_3,
\]
где $\otimes$~--- операция <<умножения>> ячеек, $\oplus$~--- операция <<сложения>> (агрегации) частичных результатов, а $\Bbbzero$~--- нейтральный элемент $\oplus$, используемый, например, для инициализации результирующей матрицы.
Само определение умножения при этом остаётся неизменным:
\[
    C[i,j] = \bigoplus_{k=0}^{n-1} A[i,k] \otimes B[k,j].
\]
Никаких дополнительных требований (ассоциативность, коммутативность, дистрибутивность) к операциям $\otimes$ и $\oplus$ не предъявляется~--- достаточно, чтобы их сигнатуры были согласованы с типами аргументов и результата%
\sidenote{Речь идёт не просто о перегрузке операторов, а о полноценных пользовательских функциях, передаваемых в качестве аргументов обобщённой процедуре умножения.}.
Так, в сообществе GraphBLAS термином <<полукольцо>> могут называть конструкцию, у которой типы аргументов умножения различны, что формально полукольцом в алгебраическом смысле не является; такая вольность оправдана практическими потребностями и устоялась в литературе%
\sidenote{
    За деталями обсуждения терминологии мы отсылаем читателя к~\cite{7761646}, а также~\cite{Garbar_Erin_Chernikov_Panfilyonok_Grigorev_2026}.
}.

При определении матрицы смежности (раздел~\ref{sec:GraphTheoryBasicDefinitions}) мы ввели множество $\Opt{L} = \{\Bbbzero\} \cup L$, где $\Bbbzero$ означает <<ребро отсутствует>>.
В языках программирования такой конструкции естественно соответствует тип \texttt{Option} (OCaml, F\#) или \texttt{Maybe} (Haskell): значение \texttt{Some x} означает <<есть ребро с меткой \texttt{x}>>, а \texttt{None}~--- <<ребра нет>>.
Таким образом, матрица смежности графа с метками типа \texttt{Lbl} имеет тип $\texttt{Matrix}\langle\texttt{Option}\langle\texttt{Lbl}\rangle\rangle$
\sidenote{%
    Можно показать, что в данном контексте $\texttt{Option}$ естественным образом образует моноид: $\texttt{None}$ играет роль $\Bbbzero$, а композиция двух значений \texttt{Some} даёт \texttt{Some} (объединение, сложение или иная операция, определённая над $L$).
}.
В реальных графах число рёбер, как правило, существенно меньше $|V|^2$, поэтому матрицы смежности разрежены.
Способы компактного хранения таких матриц подробно рассмотрены в разделе~\ref{sec:SparseMatricesAndVectors}.
Здесь же отметим ключевой для наших целей факт: разреженная матрица хранит только <<непустые>> ячейки (те, что соответствуют \texttt{Some x}), а ячейки со значением \texttt{None} восстанавливаются неявно и не хранятся.
Это порождает известную в сообществе GraphBLAS \emph{проблему явных и неявных нулей}~\sidecite{7761646}: комбинация двух не-хранимых ячеек не должна порождать новую хранимую ячейку.

Перейдём к операциям над разреженными матрицами.
\emph{Поэлементная операция} над двумя матрицами~--- это операция, которая применяется независимо к каждой паре ячеек с одинаковыми индексами.
Типизация такой операции имеет вид
\[
    \texttt{Option}\langle T_1\rangle \to \texttt{Option}\langle T_2\rangle \to \texttt{Option}\langle T_3\rangle.
\]
Заметим, что эта сигнатура допускает определение операции, возвращающей \texttt{Some x} даже тогда, когда оба аргумента равны \texttt{None}, что при работе с разреженными матрицами нежелательно%
\sidenote{Данная проблема подробно разбирается в~\cite{Garbar_Erin_Chernikov_Panfilyonok_Grigorev_2026}, где предложено решение на уровне системы типов.}.
Тем не менее, при дисциплинированном использовании эта сигнатура вполне работоспособна, и именно на её основе можно построить универсальную поэлементную операцию $\texttt{map2}$~\sidecite{7761646}:
\[
    \texttt{map2}\ (\texttt{op}: \texttt{Option}\langle T_1\rangle \to \texttt{Option}\langle T_2\rangle \to \texttt{Option}\langle T_3\rangle)\ (\texttt{m}_1, \texttt{m}_2) = \texttt{result}.
\]
Одна и та же функция $\texttt{map2}$, параметризованная разными операциями $\texttt{op}$, реализует и поэлементное сложение, и поэлементное умножение, и маскирование.

Рассмотрим конкретные примеры для целых чисел.

Операция поэлементного сложения:
\begin{align*}
    \texttt{op\_int\_add}\ (\texttt{Some}\ x)\ (\texttt{Some}\ y) &= \texttt{Some}\ (x + y), \\
    \texttt{op\_int\_add}\ (\texttt{Some}\ x)\ \texttt{None} &= \texttt{Some}\ x, \\
    \texttt{op\_int\_add}\ \texttt{None}\ (\texttt{Some}\ y) &= \texttt{Some}\ y, \\
    \texttt{op\_int\_add}\ \texttt{None}\ \texttt{None} &= \texttt{None}.
\end{align*}
Если в результате сложения получился <<честный>> ноль, можно дополнительно вернуть \texttt{None}, чтобы не сохранять его.

Операция поэлементного умножения, напротив, возвращает непустое значение только тогда, когда оба аргумента непусты:
\begin{align*}
    \texttt{op\_int\_mult}\ (\texttt{Some}\ x)\ (\texttt{Some}\ y) &= \texttt{Some}\ (x * y), \\
    \texttt{op\_int\_mult}\ \_\ \_ &= \texttt{None}.
\end{align*}

Маскирование~--- ещё один важный частный случай $\texttt{map2}$.
\emph{Структурная маска} оставляет значения первой матрицы только там, где во второй матрице есть непустые значения:
\[
    \texttt{op\_mask}\ (\texttt{Some}\ x)\ (\texttt{Some}\ y) = \texttt{Some}\ x,
\]
и \texttt{None} в остальных случаях.

\emph{Инвертированная структурная маска} действует наоборот: оставляет значения первой матрицы только там, где во второй матрице ячейка пуста.
Такая операция используется, например, в алгоритме BFS (раздел~\ref{sec:BFS}) при формировании нового фронта с учётом уже посещённых вершин:
\[
    \texttt{op\_inv\_mask}\ (\texttt{Some}\ x)\ \texttt{None} = \texttt{Some}\ x,
\]
и \texttt{None} в остальных случаях.

Наконец, можно определить \emph{предикатную маску}: оставить значение первой матрицы, если значение второй матрицы (когда оно есть) удовлетворяет заданному предикату $p$.
Структурная и инвертированная маски являются её частными случаями~--- с предикатами <<аргумент непуст>> и <<аргумент пуст>> соответственно.

Идея параметризации операций естественным образом переносится и на матричное умножение.
В терминах типов этому соответствует функция $\texttt{mxm}$:
\begin{align*}
    \texttt{mxm}\ & : \\
    & \quad \texttt{m}_1: \texttt{Matrix}\langle\texttt{Option}\langle T_1\rangle\rangle \\
    & \to \texttt{m}_2: \texttt{Matrix}\langle\texttt{Option}\langle T_2\rangle\rangle \\
    & \to \texttt{op\_mult}: \texttt{Option}\langle T_1\rangle \to \texttt{Option}\langle T_2\rangle \to \texttt{Option}\langle T_3\rangle \\
    & \to \texttt{op\_add}: \texttt{Option}\langle T_3\rangle \to \texttt{Option}\langle T_3\rangle \to \texttt{Option}\langle T_3\rangle \\
    & \to \texttt{zero}: \texttt{Option}\langle T_3\rangle \\
    & \to \texttt{result}: \texttt{Matrix}\langle\texttt{Option}\langle T_3\rangle\rangle.
\end{align*}
Здесь $\texttt{op\_mult}$ задаёт, как перемножить ячейку первой матрицы с ячейкой второй, $\texttt{op\_add}$~--- как сложить несколько частичных произведений для одной позиции результата, а $\texttt{zero}$~--- нейтральный элемент $\texttt{op\_add}$.

Вычисление происходит по формуле
\[
    \texttt{result}[i,j] = \texttt{op\_add}_{k}\ \texttt{op\_mult}\ \texttt{m}_1[i,k]\ \texttt{m}_2[k,j],
\]
где $\texttt{op\_add}_k$ обозначает итеративное применение $\texttt{op\_add}$ по $k$, начиная с $\texttt{zero}$.

Отметим важное различие с классическим определением матричного умножения над полукольцом (раздел~\ref{sec:PathAlgebra}).
В рамках описанного подхода $\texttt{op\_mult}$ и $\texttt{op\_add}$ не обязаны удовлетворять аксиомам ассоциативности, коммутативности или дистрибутивности.
На практике при решении задач анализа графов используемые операции часто не образуют полукольцо в математическом смысле, но это не мешает им быть полезными: достаточно, чтобы они были согласованы с конкретной решаемой задачей.
Когда же операции действительно образуют полукольцо, $\texttt{mxm}$ в точности совпадает с умножением матриц над этим полукольцом.

При решении некоторых задач анализа графов поэлементным операциям бывает полезно иметь доступ к координатам обрабатываемых ячеек~--- фактически, к номерам вершин.
В этом случае сигнатура операции <<умножения>> расширяется:
\[
    \otimes: (\texttt{int} \times \texttt{int} \times T_1) \times (\texttt{int} \times \texttt{int} \times T_2) \to T_3,
\]
где первые два аргумента каждого кортежа~--- номера строки и столбца обрабатываемой ячейки.
Эта возможность особенно полезна, когда результат операции зависит от того, между какими именно вершинами находится ребро.

Рассмотрим ещё один полезный на практике приём~--- выделение подмножества рёбер с использованием диагональных матриц.

Умножение матрицы смежности $M$ на диагональную матрицу $D$ справа позволяет выбрать только те рёбра, конечная вершина которых принадлежит заданному подмножеству $S \subseteq V$ (диагональ $D[i,i] = 1 \iff i \in S$).
Иными словами, $M \times D$ оставляет входящие в $S$ рёбра.

\begin{example}
    Рассмотрим граф $\mbfscrG_1$ из примера~\ref{exmpl:graph_visualization}.
    Выделим входящие рёбра для вершин $S = \{0, 2\}$ (на рисунке~\ref{fig:incomingEdges} выделены синим и красным цветами соответственно).
    Умножение справа на диагональную матрицу $D = \operatorname{diag}(1, 0, 1, 0)$ даёт:

    \begin{marginfigure}
    \begin{center}
        \begin{tikzpicture}[shorten >=1pt,auto]
  \node[state, fill=blue!20] (q_0)                      {$0$};
  \node[state] (q_1) [above right=of q_0]               {$1$};
  \node[state, fill=red!20] (q_2) [right=of q_0]        {$2$};
  \node[state] (q_3) [right=of q_2]                     {$3$};
  \path[->]
    (q_0) edge  node {} (q_1)
    (q_1) edge  node {} (q_2)
    (q_2) edge  node {} (q_0)
    (q_2) edge[bend left, above]  node {} (q_3)
    (q_3) edge[bend left, below]  node {} (q_2);
  {
    \path[->, red, thick]
    (q_1) edge node {} (q_2)
    (q_2) edge  node {} (q_0)
    (q_3) edge[bend left, below]  node {} (q_2);
  }
\end{tikzpicture}

    \end{center}
    \caption{Фильтрация входящих рёбер для вершин 0 и 2}
    \label{fig:incomingEdges}
\end{marginfigure}

    \[
        \left(
        \begin{array}{>{\columncolor{blue!20}}cc>{\columncolor{red!20}}cc}
            0 & 1 & 0 & 0 \\
            0 & 0 & 1 & 0 \\
            1 & 0 & 0 & 1 \\
            0 & 0 & 1 & 0
        \end{array}\right)
        \times
        \left(
        \begin{array}{>{\columncolor{blue!20}}cc>{\columncolor{red!20}}cc}
            1 & 0 & 0 & 0 \\
            0 & 0 & 0 & 0 \\
            0 & 0 & 1 & 0 \\
            0 & 0 & 0 & 0
        \end{array}\right)
        =
        \left(
        \begin{array}{cccc}
            0 & 0 & 0 & 0 \\
            0 & 0 & \cellcolor{red!20}{1} & 0 \\
            \cellcolor{blue!20}{1} & 0 & 0 & 0 \\
            0 & 0 & \cellcolor{red!20}{1} & 0
        \end{array}\right).
    \]
    Выделенные цветом ячейки соответствуют рёбрам: $(2,0)$~--- входящее в вершину $0$, и $(1,2)$, $(3,2)$~--- входящие в вершину $2$.
\end{example}

Умножение слева на диагональную матрицу $D \times M$, напротив, выбирает рёбра, исходящие из вершин заданного подмножества $S$:
\[
    (D \times M)[i,j] = D[i,i] \times M[i,j].
\]

\begin{example}
    Для того же графа и того же подмножества $S = \{0, 2\}$ умножение слева даёт фильтрацию исходящих рёбер (рисунок~\ref{fig:outgoingEdges}):

\begin{marginfigure}
    \begin{center}
        \input{part_01_Prep/chapter_03_GraphTheoryIntro/figures/04_AppliedAspects/outgoing_edges.tex}
    \end{center}
    \caption{Фильтрация исходящих рёбер для вершин 0 и 2 }
    \label{fig:outgoingEdges}
\end{marginfigure}


    \[
        \left(
        \begin{array}{cccc}
            \rowcolor{blue!20}
            1 & 0 & 0 & 0 \\
            0 & 0 & 0 & 0 \\
            \rowcolor{red!20}
            0 & 0 & 1 & 0 \\
            0 & 0 & 0 & 0
        \end{array}\right)
        \times
        \left(
        \begin{array}{cccc}
            \rowcolor{blue!20}
            0 & 1 & 0 & 0 \\
            0 & 0 & 1 & 0 \\
            \rowcolor{red!20}
            1 & 0 & 0 & 1 \\
            0 & 0 & 1 & 0
        \end{array}\right)
        =
        \left(
        \begin{array}{cccc}
            0 & \cellcolor{blue!20}{1} & 0 & 0 \\
            0 & 0 & 0 & 0 \\
            \cellcolor{red!20}{1} & 0 & 0 & \cellcolor{red!20}{1} \\
            0 & 0 & 0 & 0
        \end{array}\right).
    \]
    Выделены ячейки, соответствующие рёбрам: $(0,1)$~--- исходящее из вершины $0$, и $(2,0)$, $(2,3)$~--- исходящие из вершины $2$.
\end{example}

Разумеется, описанные выше типобезопасные абстракции ($\texttt{Option}$, $\texttt{map2}$, $\texttt{mxm}$) являются идеализацией.
В реальном стандарте GraphBLAS~\sidecite{7761646}~--- открытом стандарте, описывающем набор примитивов и операций разреженной линейной алгебры для построения алгоритмов анализа графов,~--- используется более <<плоское>> представление данных (без типа $\texttt{Option}$ в сигнатурах), однако ключевые идеи сохраняются: операции параметризуются пользовательскими полукольцами и моноидами, а поэлементным операциям доступны индексы обрабатываемых ячеек%
\sidenote{
    С точки зрения строгой математической терминологии именование произвольной параметризованной операции <<полукольцом>> или <<моноидом>> не вполне корректно.
    Однако идея близка, и для практических целей сообщество GraphBLAS не изобретает новых терминов, пользуясь устоявшимися названиями с некоторой вольностью.
    Детальное обсуждение этого вопроса можно найти в~\cite{7761646,Garbar_Erin_Chernikov_Panfilyonok_Grigorev_2026}.
}.
Стандарт позволяет единообразно выражать широкий класс графовых алгоритмов~--- от поиска кратчайших путей (тропическое полукольцо) до вычисления достижимости (булево полукольцо).

На практике этот стандарт реализуется в нескольких библиотеках: наиболее известная из них~--- SuiteSparse:GraphBLAS~\sidecite{Davis2018Algorithm9S}, широко используемая в высокопроизводительных графовых вычислениях.
Активно развиваются и GPGPU-реализации, такие как GraphBLAST.

Независимо от конкретной реализации, главное преимущество GraphBLAS-подхода для программиста~--- возможность выражать графовые алгоритмы через композицию матричных операций, делегируя заботу о разреженности, параллелизме и низкоуровневых оптимизациях библиотеке.
Операции $\texttt{map2}$ и $\texttt{mxm}$, параметризованные пользовательскими функциями, являются естественным обобщением этого подхода: они позволяют отказаться от жёсткого требования задавать полукольцо и вместо этого указывать ровно те операции, которые нужны для конкретной задачи, с явными гарантиями корректности на уровне типов~\sidecite{7761646}.

\input{part_01_Prep/chapter_03_GraphTheoryIntro/05_BFS}
